\IfFileExists{stacks-project.cls}{%
\documentclass{stacks-project}
}{%
\documentclass{amsart}
}

% For dealing with references we use the comment environment
\usepackage{verbatim}
\newenvironment{reference}{\comment}{\endcomment}
%\newenvironment{reference}{}{}
\newenvironment{slogan}{\comment}{\endcomment}
\newenvironment{history}{\comment}{\endcomment}

% For commutative diagrams we use Xy-pic
\usepackage[all]{xy}

% We use 2cell for 2-commutative diagrams.
\xyoption{2cell}
\UseAllTwocells

% We use multicol for the list of chapters between chapters
\usepackage{multicol}

% This is generally recommended for better output
\usepackage{lmodern}
\usepackage[T1]{fontenc}

% For cross-file-references
\usepackage{xr-hyper}

% Package for hypertext links:
\usepackage{hyperref}

% For any local file, say "hello.tex" you want to link to please
% use \externaldocument[hello-]{hello}
\externaldocument[introduction-]{introduction}
\externaldocument[conventions-]{conventions}
\externaldocument[sets-]{sets}
\externaldocument[categories-]{categories}
\externaldocument[topology-]{topology}
\externaldocument[sheaves-]{sheaves}
\externaldocument[sites-]{sites}
\externaldocument[stacks-]{stacks}
\externaldocument[fields-]{fields}
\externaldocument[algebra-]{algebra}
\externaldocument[brauer-]{brauer}
\externaldocument[homology-]{homology}
\externaldocument[derived-]{derived}
\externaldocument[simplicial-]{simplicial}
\externaldocument[more-algebra-]{more-algebra}
\externaldocument[smoothing-]{smoothing}
\externaldocument[modules-]{modules}
\externaldocument[sites-modules-]{sites-modules}
\externaldocument[injectives-]{injectives}
\externaldocument[cohomology-]{cohomology}
\externaldocument[sites-cohomology-]{sites-cohomology}
\externaldocument[dga-]{dga}
\externaldocument[dpa-]{dpa}
\externaldocument[sdga-]{sdga}
\externaldocument[hypercovering-]{hypercovering}
\externaldocument[schemes-]{schemes}
\externaldocument[constructions-]{constructions}
\externaldocument[properties-]{properties}
\externaldocument[morphisms-]{morphisms}
\externaldocument[coherent-]{coherent}
\externaldocument[divisors-]{divisors}
\externaldocument[limits-]{limits}
\externaldocument[varieties-]{varieties}
\externaldocument[topologies-]{topologies}
\externaldocument[descent-]{descent}
\externaldocument[perfect-]{perfect}
\externaldocument[more-morphisms-]{more-morphisms}
\externaldocument[flat-]{flat}
\externaldocument[groupoids-]{groupoids}
\externaldocument[more-groupoids-]{more-groupoids}
\externaldocument[etale-]{etale}
\externaldocument[chow-]{chow}
\externaldocument[intersection-]{intersection}
\externaldocument[pic-]{pic}
\externaldocument[weil-]{weil}
\externaldocument[adequate-]{adequate}
\externaldocument[dualizing-]{dualizing}
\externaldocument[duality-]{duality}
\externaldocument[discriminant-]{discriminant}
\externaldocument[derham-]{derham}
\externaldocument[local-cohomology-]{local-cohomology}
\externaldocument[algebraization-]{algebraization}
\externaldocument[curves-]{curves}
\externaldocument[resolve-]{resolve}
\externaldocument[models-]{models}
\externaldocument[functors-]{functors}
\externaldocument[equiv-]{equiv}
\externaldocument[pione-]{pione}
\externaldocument[etale-cohomology-]{etale-cohomology}
\externaldocument[proetale-]{proetale}
\externaldocument[relative-cycles-]{relative-cycles}
\externaldocument[more-etale-]{more-etale}
\externaldocument[trace-]{trace}
\externaldocument[crystalline-]{crystalline}
\externaldocument[spaces-]{spaces}
\externaldocument[spaces-properties-]{spaces-properties}
\externaldocument[spaces-morphisms-]{spaces-morphisms}
\externaldocument[decent-spaces-]{decent-spaces}
\externaldocument[spaces-cohomology-]{spaces-cohomology}
\externaldocument[spaces-limits-]{spaces-limits}
\externaldocument[spaces-divisors-]{spaces-divisors}
\externaldocument[spaces-over-fields-]{spaces-over-fields}
\externaldocument[spaces-topologies-]{spaces-topologies}
\externaldocument[spaces-descent-]{spaces-descent}
\externaldocument[spaces-perfect-]{spaces-perfect}
\externaldocument[spaces-more-morphisms-]{spaces-more-morphisms}
\externaldocument[spaces-flat-]{spaces-flat}
\externaldocument[spaces-groupoids-]{spaces-groupoids}
\externaldocument[spaces-more-groupoids-]{spaces-more-groupoids}
\externaldocument[bootstrap-]{bootstrap}
\externaldocument[spaces-pushouts-]{spaces-pushouts}
\externaldocument[spaces-chow-]{spaces-chow}
\externaldocument[groupoids-quotients-]{groupoids-quotients}
\externaldocument[spaces-more-cohomology-]{spaces-more-cohomology}
\externaldocument[spaces-simplicial-]{spaces-simplicial}
\externaldocument[spaces-duality-]{spaces-duality}
\externaldocument[formal-spaces-]{formal-spaces}
\externaldocument[restricted-]{restricted}
\externaldocument[spaces-resolve-]{spaces-resolve}
\externaldocument[formal-defos-]{formal-defos}
\externaldocument[defos-]{defos}
\externaldocument[cotangent-]{cotangent}
\externaldocument[examples-defos-]{examples-defos}
\externaldocument[algebraic-]{algebraic}
\externaldocument[examples-stacks-]{examples-stacks}
\externaldocument[stacks-sheaves-]{stacks-sheaves}
\externaldocument[criteria-]{criteria}
\externaldocument[artin-]{artin}
\externaldocument[quot-]{quot}
\externaldocument[stacks-properties-]{stacks-properties}
\externaldocument[stacks-morphisms-]{stacks-morphisms}
\externaldocument[stacks-limits-]{stacks-limits}
\externaldocument[stacks-cohomology-]{stacks-cohomology}
\externaldocument[stacks-perfect-]{stacks-perfect}
\externaldocument[stacks-introduction-]{stacks-introduction}
\externaldocument[stacks-more-morphisms-]{stacks-more-morphisms}
\externaldocument[stacks-geometry-]{stacks-geometry}
\externaldocument[moduli-]{moduli}
\externaldocument[moduli-curves-]{moduli-curves}
\externaldocument[examples-]{examples}
\externaldocument[exercises-]{exercises}
\externaldocument[guide-]{guide}
\externaldocument[desirables-]{desirables}
\externaldocument[coding-]{coding}
\externaldocument[obsolete-]{obsolete}
\externaldocument[fdl-]{fdl}
\externaldocument[index-]{index}

% Theorem environments.
%
\theoremstyle{plain}
\newtheorem{theorem}[subsection]{Theorem}
\newtheorem{proposition}[subsection]{Proposition}
\newtheorem{lemma}[subsection]{Lemma}

\theoremstyle{definition}
\newtheorem{definition}[subsection]{Definition}
\newtheorem{example}[subsection]{Example}
\newtheorem{exercise}[subsection]{Exercise}
\newtheorem{situation}[subsection]{Situation}

\theoremstyle{remark}
\newtheorem{remark}[subsection]{Remark}
\newtheorem{remarks}[subsection]{Remarks}

\numberwithin{equation}{subsection}

% Macros
%
\def\lim{\mathop{\mathrm{lim}}\nolimits}
\def\colim{\mathop{\mathrm{colim}}\nolimits}
\def\Spec{\mathop{\mathrm{Spec}}}
\def\Hom{\mathop{\mathrm{Hom}}\nolimits}
\def\Ext{\mathop{\mathrm{Ext}}\nolimits}
\def\SheafHom{\mathop{\mathcal{H}\!\mathit{om}}\nolimits}
\def\SheafExt{\mathop{\mathcal{E}\!\mathit{xt}}\nolimits}
\def\Sch{\mathit{Sch}}
\def\Mor{\mathop{\mathrm{Mor}}\nolimits}
\def\Ob{\mathop{\mathrm{Ob}}\nolimits}
\def\Sh{\mathop{\mathit{Sh}}\nolimits}
\def\NL{\mathop{N\!L}\nolimits}
\def\CH{\mathop{\mathrm{CH}}\nolimits}
\def\proetale{{pro\text{-}\acute{e}tale}}
\def\etale{{\acute{e}tale}}
\def\QCoh{\mathit{QCoh}}
\def\Ker{\mathop{\mathrm{Ker}}}
\def\Im{\mathop{\mathrm{Im}}}
\def\Coker{\mathop{\mathrm{Coker}}}
\def\Coim{\mathop{\mathrm{Coim}}}

% Boxtimes
%
\DeclareMathSymbol{\boxtimes}{\mathbin}{AMSa}{"02}

%
% Macros for moduli stacks/spaces
%
\def\QCohstack{\mathcal{QC}\!\mathit{oh}}
\def\Cohstack{\mathcal{C}\!\mathit{oh}}
\def\Spacesstack{\mathcal{S}\!\mathit{paces}}
\def\Quotfunctor{\mathrm{Quot}}
\def\Hilbfunctor{\mathrm{Hilb}}
\def\Curvesstack{\mathcal{C}\!\mathit{urves}}
\def\Polarizedstack{\mathcal{P}\!\mathit{olarized}}
\def\Complexesstack{\mathcal{C}\!\mathit{omplexes}}
% \Pic is the operator that assigns to X its picard group, usage \Pic(X)
% \Picardstack_{X/B} denotes the Picard stack of X over B
% \Picardfunctor_{X/B} denotes the Picard functor of X over B
\def\Pic{\mathop{\mathrm{Pic}}\nolimits}
\def\Picardstack{\mathcal{P}\!\mathit{ic}}
\def\Picardfunctor{\mathrm{Pic}}
\def\Deformationcategory{\mathcal{D}\!\mathit{ef}}
\usepackage{cleveref}
\crefname{originalreference}{original reference}{original references}
\Crefname{originalreference}{Original reference}{Original references}
\makeatletter
\newlabel{schemes-lemma-glue-functors}{{01JJ}{source}{}{sourcepointer-1}{}}
\newlabel{schemes-lemma-glue-functors@cref}{{[originalreference][0][]01JJ}{[1][0][]source}{}{}{}}
\makeatother
\title{Glueing schemes over a basis from compatible local data}
\author{The Stacks Project Authors}
\date{stacks-project:constructions.tex:local-snapshot-20261012; GNU FDL 1.2 or later}
\begin{document}
\maketitle
\noindent\textit{Editorial note.} One '(easy) verification' omitted: the constructed system (g, \{h\_U\}) lies in F\_W(T) and maps to alpha under the displayed arrow; the construction of the h\_U themselves is fully displayed via conditions (a) and (b).\par\medskip
\noindent\textit{Allowed standard prerequisites (editorial).}\begin{itemize}
\item Glue-functors criterion. schemes-lemma-glue-functors (cross-chapter, named). Provides the four-step representability template.
\item Glueing of sheaves. sheaves-section-glueing-sheaves (cross-chapter, named). Underlies the sheaf property of the constructed functor.
\end{itemize}\medskip
\ifdefined\hypertarget\hypertarget{sourcepointer-1}{}\fi
\noindent\textit{Source reference.} Original reference 01JJ: focused-seeds/geometry/S9-cons-relative-glueing/ref-tags.txt, source lines 5--5. Navigation pointer only: the counted proof cites this statement by \textbackslash{}ref as a standing input; the deduction the counted proof has to make is carried out inside the counted span, and no part of that deduction is delegated to this pointer.\par\medskip
\begin{lemma}
\label{lemma-relative-glueing}
Let $S$ be a scheme.
Let $\mathcal{B}$ be a basis for the topology of $S$.
Suppose given the following data:
\begin{enumerate}
\item For every $U \in \mathcal{B}$ a scheme $f_U : X_U \to U$ over $U$.
\item For $U, V \in \mathcal{B}$ with $V \subset U$ a morphism
$\rho^U_V : X_V \to X_U$ over $U$.
\end{enumerate}
Assume that
\begin{enumerate}
\item[(a)] each $\rho^U_V$ induces an isomorphism
$X_V \to f_U^{-1}(V)$ of schemes over $V$,
\item[(b)] whenever $W, V, U \in \mathcal{B}$, with
$W \subset V \subset U$ we have $\rho^U_W = \rho^U_V \circ \rho ^V_W$.
\end{enumerate}
Then there exists a morphism $f : X \to S$ of schemes
and isomorphisms $i_U : f^{-1}(U) \to X_U$ over $U \in \mathcal{B}$
such that for $V, U \in \mathcal{B}$ with $V \subset U$ the composition
$$
\xymatrix{
X_V \ar[r]^{i_V^{-1}} &
f^{-1}(V) \ar[rr]^{inclusion} & &
f^{-1}(U) \ar[r]^{i_U} &
X_U
}
$$
is the morphism $\rho^U_V$. Moreover $X$ is unique up to
unique isomorphism over $S$.
\end{lemma}

\begin{proof}
To prove this we will use Schemes, Lemma \ref{schemes-lemma-glue-functors}.
First we define a contravariant functor $F$ from the category of schemes
to the category of sets. Namely, for a scheme $T$ we set
$$
F(T) =
\left\{
\begin{matrix}
(g, \{h_U\}_{U \in \mathcal{B}}),
\ g : T \to S, \ h_U : g^{-1}(U) \to X_U, \\
f_U \circ h_U = g|_{g^{-1}(U)},
\ h_U|_{g^{-1}(V)} = \rho^U_V \circ h_V
\ \forall\ V, U \in \mathcal{B}, V \subset U
\end{matrix}
\right\}.
$$
The restriction mapping $F(T) \to F(T')$ given a morphism
$T' \to T$ is just gotten by composition.
For any $W \in \mathcal{B}$ we consider the subfunctor
$F_W \subset F$ consisting of those systems $(g, \{h_U\})$
such that $g(T) \subset W$.

\medskip\noindent
First we show $F$ satisfies the sheaf property for the Zariski topology.
Suppose that $T$ is a scheme, $T = \bigcup V_i$ is an open covering,
and $\xi_i \in F(V_i)$ is an element such that
$\xi_i|_{V_i \cap V_j} = \xi_j|_{V_i \cap V_j}$.
Say $\xi_i = (g_i, \{h_{i, U}\})$. Then we immediately see that
the morphisms $g_i$ glue to a unique global morphism
$g : T \to S$. Moreover, it is clear that
$g^{-1}(U) = \bigcup g_i^{-1}(U)$. Hence the morphisms
$h_{i, U} : g_i^{-1}(U) \to X_U$ glue to a unique morphism
$h_U : g^{-1}(U) \to X_U$. It is easy to verify that the system
$(g, \{h_U\})$ is an element of $F(T)$. Hence $F$ satisfies the
sheaf property for the Zariski topology.

\medskip\noindent
Next we verify that each $F_W$, $W \in \mathcal{B}$ is representable.
Namely, we claim that the transformation of functors
$$
F_W \longrightarrow \Mor(-, X_W), \ (g, \{h_U\}) \longmapsto h_W
$$
is an isomorphism. To see this suppose that $T$ is a scheme and
$\alpha : T \to X_W$ is a morphism. Set $g = f_W \circ \alpha$.
For any $U \in \mathcal{B}$ such that $U \subset W$ we can
define $h_U : g^{-1}(U) \to X_U$ be the composition
$(\rho^W_U)^{-1} \circ \alpha|_{g^{-1}(U)}$. This works because
the image $\alpha(g^{-1}(U))$ is contained in $f_W^{-1}(U)$ and
condition (a) of the lemma. It is clear that
$f_U \circ h_U = g|_{g^{-1}(U)}$ for such a $U$.
Moreover, if also $V \in \mathcal{B}$ and $V \subset U \subset W$,
then $\rho^U_V \circ h_V = h_U|_{g^{-1}(V)}$ by property (b)
of the lemma. We still have to define $h_U$ for an arbitrary
element $U \in \mathcal{B}$. Since $\mathcal{B}$ is a basis for
the topology on $S$ we can find an open covering
$U \cap W = \bigcup U_i$ with $U_i \in \mathcal{B}$. Since $g$ maps into $W$
we have
$g^{-1}(U) = g^{-1}(U \cap W) = \bigcup g^{-1}(U_i)$.
Consider the morphisms
$h_i = \rho^U_{U_i} \circ h_{U_i} : g^{-1}(U_i) \to X_U$.
It is a simple matter to use condition (b) of the lemma
to prove that
$h_i|_{g^{-1}(U_i) \cap g^{-1}(U_j)} = h_j|_{g^{-1}(U_i) \cap g^{-1}(U_j)}$.
Hence these morphisms glue to give the desired morphism
$h_U : g^{-1}(U) \to X_U$. We omit the (easy) verification that
the system $(g, \{h_U\})$ is an element of $F_W(T)$ which
maps to $\alpha$ under the displayed arrow above.

\medskip\noindent
Next, we verify each $F_W \subset F$ is representable by open immersions.
This is clear from the definitions.

\medskip\noindent
Finally we have to verify
the collection $(F_W)_{W \in \mathcal{B}}$ covers $F$.
This is clear by construction and the fact that $\mathcal{B}$ is
a basis for the topology of $S$.

\medskip\noindent
Let $X$ be a scheme representing the functor $F$.
Let $(f, \{i_U\}) \in F(X)$ be a ``universal family''.
Since each $F_W$ is representable by $X_W$ (via the morphism of functors
displayed above) we see that $i_W : f^{-1}(W) \to X_W$
is an isomorphism as desired. The lemma is proved.
\end{proof}

\end{document}
